\scanpage{90}
\mspage{\handwritten{57}}
limit ourselves to the case $g=2$. We are interested in $M(\Gamma_{2}(\unsure{\lambda|l}))$ for $\lambda=1,2$, because we know $S(\Gamma_{2}(\lambda))$ explicitly for the same values of $\lambda$ (8). We keep in mind that they are defined as quotient varieties and not as monoidal transforms. We take a point $t_{o}$ of $\underline{S}_{1}$. Since $t_{o}$ is not a fixed point of any element of $\Gamma_{1}(2)$ different from $\pm1_{2}$, the analytic local ring $\unsure{O|0}$ of $M(\Gamma_{2}(2))$ at the image point of $t_{o}$ consists of convergent series of the form
\[
\sum_{k=\unsure{o|0}}^{\infty}\theta_{k}(t,z)e(kw/2)\,.
\]
The coefficient $\theta_{k}(t,z)$ is holomorphic in $V\handwritten{\times}\mathbb{C}$ for some $V$ and $z\longrightarrow\theta_{k}(t,2z)$ is a theta-function of order $\handwritten{4}k$, which is even in the sense it is invariant under $z\longrightarrow-z$. This follows from Theorem 1. Moreover, the ideal $I$ consists of convergent series which start from $k=1$. Furthermore, there exists a base $(\theta_{\unsure{l|1}}(t_{o},z))$ over $\mathbb{C}$ of the vector space of even theta-functions of order four such that each $\theta_{l}(t,z)$ is holomorphic in some $V\handwritten{\times}\mathbb{C}$ and such that the correspondence $z\longrightarrow(\theta_{l}(t,z))$ gives rise to a complex-analytic embedding of the quotient variety of $T_{1}(t)$ by its involution $z\longrightarrow-z$, which is $P_{1}(\mathbb{C})$, as a plane curve of order two. Therefore, if we take the monoidal transform of $S(\Gamma_{2}(2))$ along its boundary in the neighborhood of the image point of $t_{o}$, we get the quotient variety of the monoidal transform of $S(\Gamma_{2}(\lambda))$, for any even $\lambda$ at least equal to 3, along its boundary in the neighborhood of the image point of $t_{o}$. In the same way, we can conclude that the morphism $M(\Gamma_{2}(1))\longrightarrow S(\Gamma_{2}(1))$ is a monoidal transformation in the neighborhood of the image point of $t_{o}$ provided $t_{o}$ is not a fixed point of any element of $\Gamma_{1}(1)$ different

% archive sheet number: 89

